% ── Sección: Instalación ──────────────────────────────────────────────────────
\section{Instalación}

% ── El problema ───────────────────────────────────────────────────────────────
\begin{frame}{El problema con LaTeX tradicional}
  \begin{columns}[T]
    \begin{column}{0.5\textwidth}
      \textbf{Lo que normalmente pasa:}
      \begin{itemize}
        \item \texttt{sudo apt install texlive-full} \\ {\small (20 min, 4 GB)}
        \item Conflictos de versiones entre paquetes
        \item Funciona en mi maquina, no en la tuya
        \item Reinstalar todo en cada equipo nuevo
      \end{itemize}
    \end{column}
    \begin{column}{0.5\textwidth}
      \textbf{Con pptex:}
      \begin{itemize}
        \item Solo necesitas \alert{Docker} y \alert{Git}
        \item Mismo entorno en cualquier OS
        \item Un solo comando para empezar
        \item Cero conflictos, cero configuracion
      \end{itemize}
    \end{column}
  \end{columns}
\end{frame}

% ── Requisitos ────────────────────────────────────────────────────────────────
\begin{frame}[fragile]{Requisitos: solo dos cosas}
  \begin{block}{1. Docker Desktop}
    Descarga desde \texttt{docker.com/products/docker-desktop} \\
    Disponible para macOS, Linux y Windows.
  \end{block}

  \vspace{0.8em}

  \begin{block}{2. Git}
    Ya lo tienes instalado. Si no: \texttt{git-scm.com}
  \end{block}

  \vspace{0.8em}
  \begin{alertblock}{Eso es todo}
    No necesitas instalar TeX Live, MiKTeX, ni ningun paquete LaTeX.
    Todo corre dentro del contenedor Docker.
  \end{alertblock}
\end{frame}

% ── Pasos ─────────────────────────────────────────────────────────────────────
\begin{frame}[fragile]{Instalacion en 3 comandos}

  \textbf{Paso 1 — Clonar el repositorio:}
  \begin{lstlisting}[language=bash]
git clone https://github.com/Wilmar3752/pptex.git
cd pptex
chmod +x pptex       # macOS / Linux
  \end{lstlisting}

  \textbf{Paso 2 — Construir la imagen Docker:}
  \begin{lstlisting}[language=bash]
./pptex build
  \end{lstlisting}

  \begin{exampleblock}{Solo una vez}
    La imagen se descarga y construye una sola vez (\(\sim\)4 GB).
    Despues todo es instantaneo.
  \end{exampleblock}
\end{frame}

% ── Compatibilidad ────────────────────────────────────────────────────────────
\begin{frame}{Compatibilidad por sistema operativo}
  \begin{table}
    \centering
    \begin{tabular}{llc}
      \toprule
      Sistema operativo        & Terminal recomendada & Funciona \\
      \midrule
      macOS                    & Terminal / iTerm2    & \textcolor{green!60!black}{\textbf{Si}} \\
      Linux                    & Cualquiera           & \textcolor{green!60!black}{\textbf{Si}} \\
      Windows — WSL2           & Terminal WSL         & \textcolor{green!60!black}{\textbf{Si}} \\
      Windows — Git Bash       & Git Bash             & \textcolor{green!60!black}{\textbf{Si}} \\
      Windows — PowerShell/CMD & —                    & \textcolor{red}{\textbf{No}} \\
      \bottomrule
    \end{tabular}
  \end{table}

  \vspace{0.5em}
  \begin{block}{Windows: recomendacion}
    Usar \textbf{WSL2} — la experiencia es identica a Linux y Docker
    Desktop se integra nativamente.
  \end{block}
\end{frame}
